% !TeX root = ../../Book1.tex
% 纲要标题：### 1.2 半群与幺半群
% 纲要位置：计划纲要.md，第 139 行。

\section{半群与幺半群}
\label{b1:sec:0102}

% 纲要内容（仅作写作计划，尚非教材正文）：
%
% * 半群、幺半群及同态
% * 自由幺半群与字的拼接
% * 变换幺半群、消去律与群化的动机
%

\subsection{半群、幺半群及同态}
\label{b1:subsec:0102:01}

把结合律作为共同条件，便能同时讨论字的拼接、映射复合和矩阵乘法。由于这些运算通常不交换，本节使用乘法记号 \(ab\)，但不把因子顺序视为可以任意更换。

\begin{definition}\label{b1:def:monoid}
带有结合的二元运算的非空集合称为\term[banqun]{半群}[semigroup]。具有单位元的半群称为\term[yaobanqun]{幺半群}[monoid]，其单位元一般记为 \(1\)。幺半群 \(M,N\) 之间的\term[tongtai]{同态}[homomorphism]是满足
\[
f(ab)=f(a)f(b),\qquad f(1_M)=1_N
\]
的映射 \(f:M\rightarrow N\)。半群同态只要求乘法等式。
\end{definition}

单位保持条件不能从乘法保持条件自动得到。例如从单元素幺半群到整数乘法幺半群、把唯一元素送到 \(0\) 的映射保持乘法，却不保持单位。今后“幺半群同态”总包含保持单位的要求；这使空乘积与非空乘积具有相同的映射规则。

\ProseParagraphBreak
\((\mathbb N,+)\) 是以 \(0\) 为单位元的交换幺半群；正整数加法构成半群，却没有单位元。若 \(M\) 是幺半群，递归定义 \(a^0=1\)、\(a^{n+1}=a^na\)。结合律给出 \(a^{m+n}=a^ma^n\)：固定 \(m\) 对 \(n\) 归纳即可。于是任取 \(a\in M\)，映射 \(n\mapsto a^n\) 都是从 \((\mathbb N,+)\) 到 \(M\) 的同态。

\begin{proposition}\label{b1:prop:monoid-maps}
幺半群同态的复合仍是同态；恒等映射是同态。若同态 \(f:M\rightarrow N\) 是双射，则逆映射 \(f^{-1}\) 也保持乘法和单位。
\end{proposition}
\begin{proof}
对于同态 \(f:M\rightarrow N\)、\(g:N\rightarrow P\)，有
\((gf)(ab)=g(f(a)f(b))=(gf)(a)(gf)(b)\)，并且
\((gf)(1_M)=1_P\)。恒等映射的结论直接由定义得出。若 \(f\) 双射，写 \(x=f(a),y=f(b)\)，则
\(f^{-1}(xy)=f^{-1}(f(ab))=ab=f^{-1}(x)f^{-1}(y)\)，且
\(f^{-1}(1_N)=1_M\)。
\end{proof}

这样的双射同态称为\term[tonggou]{同构}[isomorphism]。它说明两个幺半群的乘法结构相同，只是元素的名称不同。若只比较底集合的大小，就会丢掉乘法信息；后面将反复用元素阶、交换性等性质区别同样大小的结构。

\subsection{自由幺半群与字的拼接}
\label{b1:subsec:0102:02}

给定集合 \(X\)，把它的元素看成字母。长度为 \(n\) 的\term[zi]{字}[word]是有序序列 \((x_1,\ldots,x_n)\)，通常写成 \(x_1\cdots x_n\)。长度零的字只有一个，称为空字，记为 \(\varepsilon\)。不同长度的字始终视为不同对象，即使 \(X\) 自身含有数或序列。

\begin{definition}\label{b1:def:free-monoid}
集合 \(X\) 上的\term[ziyou-yaobanqun]{自由幺半群}[free monoid]是所有有限字的集合
\[
X^*\coloneqq\coprod_{n\geq0}X^n,
\]
运算为前后拼接，单位元为 \(\varepsilon\)。记 \(\iota:X\rightarrow X^*\) 为把字母送到长度一的字的映射。
\end{definition}
\symidx{\(X^*\)：自由幺半群}

符号 \(\coprod\) 表示不交并，即保留每个字的长度标记。拼接后各位置的字母就是先列出第一个字，再列出第二个字。因此两次拼接的结果不依赖括号；空字不增加任何位置，确为单位元。例如 \(X=\{a,b\}\) 时，\((ab)(baa)=abbaa\)，而 \((baa)(ab)=baaab\)。所以当至少有两个字母时，自由幺半群并不交换。单字母的情形只有 \(\varepsilon,a,a^2,\ldots\)，按长度同构于 \((\mathbb N,+)\)。

\begin{theorem}[字母映射的延拓]\label{b1:thm:free-monoid}
任给幺半群 \(M\) 和集合映射 \(u:X\rightarrow M\)，存在唯一幺半群同态
\(\widehat u:X^*\rightarrow M\) 满足 \(\widehat u\iota=u\)。它由
\[
\widehat u(\varepsilon)=1_M,\qquad
\widehat u(x_1\cdots x_n)=u(x_1)\cdots u(x_n)
\]
给出。
\end{theorem}
\begin{proof}
一个非空字有唯一的字母序列，右端乘积由结合律不依赖括号，所以公式定义了映射。两个字拼接时，右端乘积恰是两个原乘积的乘积，故它保持乘法；空字的指定值保证保持单位。每个非空字都是长度一的字的乘积，任何满足要求的同态都只能取上述值；在空字上也只能取单位元，所以延拓唯一。
\end{proof}

例如给 \(a,b\) 指定两个同阶方阵 \(A,B\)，每个字便成为按既定次序进行的矩阵乘积。可能有不同字取同一矩阵，但在 \(X^*\) 中它们仍是不同的字。“自由”所排除的是额外规定的字之间的等式，并不要求每次代入后的取值都不同。

\subsection{变换幺半群、消去律与群化的动机}
\label{b1:subsec:0102:03}

集合 \(X\) 的所有自映射在复合下构成\term[bianhuan-yaobanqun]{变换幺半群}[transformation monoid]，记为 \(\operatorname{End}_{\mathrm{Set}}(X)\)，单位元为恒等映射。前节的移位映射说明，其中可有左逆而无右逆。由自映射到置换的限制，正是从一般操作转向可逆操作。

\begin{definition}\label{b1:def:cancellation}
幺半群满足\term[xiaoqulv]{消去律}[cancellation law]，是指
\(ax=ay\Rightarrow x=y\) 与 \(xa=ya\Rightarrow x=y\) 都成立。
\end{definition}

自由幺半群满足消去律，因为等字有相同的逐位字母，去掉相同前缀或后缀仍得等字。然而非空字没有逆元：拼接后的长度是两字长度之和，不可能成为零。因此消去律本身不等于可逆性。

\ProseParagraphBreak
对交换幺半群，可以像从自然数构造整数那样加入形式差。以下给出消去律充分发挥作用的一种情形。为了使公式清晰，暂以加法写 \(M\)，单位元记为 \(0\)。

\begin{proposition}[交换消去幺半群的群化]\label{b1:prop:commutative-completion}
设 \(M\) 是满足消去律的交换幺半群。在 \(M\times M\) 上规定
\[
(a,b)\sim(c,d)\quad\Longleftrightarrow\quad a+d=c+b.
\]
其等价类集合 \(K(M)\) 在
\([(a,b)]+[(c,d)]=[(a+c,b+d)]\) 下是交换幺半群，每个元素都有逆元 \([(b,a)]\)。映射
\(j(a)=[(a,0)]\) 是单射同态。
\end{proposition}
\begin{proof}
自反性和对称性直接成立。若 \(a+d=c+b\) 且 \(c+f=e+d\)，相加后消去 \(c+d\)，得到 \(a+f=e+b\)，故有传递性。若两对代表元分别等价，则把相应两条等式相加可知它们的逐项和等价，所以加法良定义。结合律与交换律来自 \(M\)，零元为 \([(0,0)]\)，而
\([(a+b,b+a)]=[(0,0)]\)，给出所述逆元。最后，\(j(a)=j(c)\) 当且仅当 \(a=c\)，并且 \(j\) 保持加法与零元。
\end{proof}

这里“等价类”的一般性质将在下一节系统说明；上述证明已经逐项给出了本构造需要的检验。对 \(M=\mathbb N\)，映射 \([(a,b)]\mapsto a-b\) 是到 \(\mathbb Z\) 的双射并保持加法。一般情形也有同样的映射原则：若 \(u:M\rightarrow A\) 是到一个每个元素都可逆的交换幺半群的同态，则
\([(a,b)]\mapsto u(a)-u(b)\) 是唯一延拓 \(u\) 的同态。良定义来自
\(a+d=c+b\) 蕴含 \(u(a)+u(d)=u(c)+u(b)\)；唯一性来自
\([(a,b)]=j(a)-j(b)\)。这说明新加的形式差除了使元素可逆，没有强加不必要的关系。

\ProseParagraphBreak
交换性在这里允许整理求和，消去律则用于等价关系的传递性和嵌入判断。非交换情形不能直接把同一公式中的加号换成乘号；一般半群嵌入群的问题有额外障碍，本节不作这种推广。
